% =============================================================================
%  CAPÍTULO 8: DISSERTAÇÃO, PRODUÇÃO CIENTÍFICA QUALIS A1 E DEFESA
%  Níveis: 0 (zero) ao PhD
%  Autor: Marcelo Claro Laranjeira
% =============================================================================
\chapter{Dissertação, Produção Científica Qualis A1 e Defesa Perante Banca Acadêmica}
\label{cap:dissertacao-banca}

\paragraph{Corando a jornada acadêmica} A produção acadêmica de alto nível --- representada por dissertações de mestrado,
teses de doutorado e artigos em periódicos \qualisA{} --- constitui o ápice da
validação científica no ecossistema de pesquisa brasileiro. Este capítulo
apresenta a metodologia completa para a produção, validação e defesa de trabalhos
acadêmicos utilizando o \opencodeeco{}, desde a concepção do tema até a arguição
perante a banca examinadora.

O capítulo está organizado em dez seções progressivas, que acompanham o
pesquisador do nível zero (compreensão do que é uma dissertação) ao nível PhD
(simulação de banca, auditoria Qualis A1 e produção de artigos de alto impacto).
A Tabela~\ref{tab:conteudo-capitulo8} apresenta a estrutura completa.

\begin{table}[htb]
\centering
    \small
\caption{Estrutura do Capítulo 8}
\label{tab:conteudo-capitulo8}
    \resizebox{\textwidth}{!}{
\begin{tabular}{clcc}
\toprule
\textbf{Seção} & \textbf{Tópico} & \textbf{Nível} & \textbf{Páginas} \\
\midrule
8.1 & Introdução à Produção Acadêmica com Agentes & \nivelzero & 4 \\
8.2 & Metodologia PPGTE/UFC: Estrutura da Dissertação & \nivelavancado & 10 \\
8.3 & Protocolo de Anonimato para Avaliação Cega & \nivelavancado & 6 \\
8.4 & Simulação de Banca com Agent-Forum & \nivelphd & 12 \\
8.5 & PhD Auditor: Nash, Cohen, Bonferroni, Qualis & \nivelphd & 10 \\
8.6 & AUTO\_SCORE\_QUALIS: Sistema Automático de Pontuação & \nivelavancado & 8 \\
8.7 & Iterative Correction Loop: Ciclo de Refinamento & \nivelavancado & 8 \\
8.8 & Produção de Artigos Qualis A1 & \nivelphd & 10 \\
8.9 & Roteiro do Nível Zero ao PhD & \text{Todos} & 6 \\
8.10 & Conclusão e Perspectivas Futuras & \nivelphd & 6 \\
\bottomrule
\end{tabular}
    }
\end{table}

Ao final deste capítulo, o leitor será capaz de:
\begin{itemize}
\item Compreender o percurso completo de produção acadêmica, do tema à defesa;
\item Estruturar uma dissertação nos moldes do PPGTE/UFC com template \texttt{abntex2};
\item Aplicar o protocolo de anonimato para avaliação cega;
\item Simular uma banca examinadora com agentes multiagente;
\item Utilizar o PhD Auditor para validação estatística e Qualis A1;
\item Executar o ciclo iterativo de correção até pontuação $\geq 95$;
\item Produzir artigos científicos de alto impacto com o pipeline SEEKER-MASWOS;
\item Navegar do nível zero ao PhD com um roteiro estruturado.
\end{itemize}

% =============================================================================
%  SEÇÃO 8.1: INTRODUÇÃO À PRODUÇÃO ACADÊMICA COM AGENTES
%  Nível: 0 (zero) / Básico (~4 laudas)
% =============================================================================
\section{Introdução à Produção Acadêmica com Agentes}
\label{sec:introducao-producao-academica}
\nivelzero

\paragraph{Entendendo o que significa produzir ciência com agentes} A produção acadêmica é o processo pelo qual o conhecimento científico é
sistematizado, documentado e submetido ao escrutínio da comunidade de
pesquisadores \cite{popper2005logica}. No contexto brasileiro, a dissertação de
mestrado e o artigo em periódico \qualisA{} representam os veículos mais
prestigiados de comunicação científica.

\subsection{O que é uma Dissertação?}

A dissertação de mestrado é um trabalho acadêmico que demonstra a capacidade do
candidato de realizar pesquisa científica de forma sistemática. Diferentemente de
uma tese de doutorado, que exige contribuição original significativa ao avanço do
conhecimento, a dissertação requer \destaque{domínio da metodologia científica} e
\destaque{capacidade de análise crítica} da literatura existente
\cite{abnt2023nbr14724}.

Os elementos essenciais de uma dissertação incluem:
\begin{enumerate}
\item \destaque{Problema de pesquisa}: questão central que motiva a investigação;
\item \destaque{Referencial teórico}: base conceitual que fundamenta a análise;
\item \destaque{Metodologia}: procedimentos sistemáticos para coleta e análise de dados;
\item \destaque{Resultados}: evidências obtidas através da aplicação da metodologia;
\item \destaque{Discussão}: interpretação dos resultados à luz da literatura;
\item \destaque{Conclusão}: síntese das contribuições e limitações do estudo.
\end{enumerate}

\subsection{O que é um Artigo Qualis A1?}

O \qualisA{} é o estrato mais elevado do sistema de classificação de periódicos
da CAPES (Coordenação de Aperfeiçoamento de Pessoal de Nível Superior). Periódicos
classificados como A1 representam o percentual superior (aproximadamente 25\%) dos
veículos de maior impacto e rigor científico em cada área do conhecimento
\cite{laranjeira2026opencode}.

Publicar em um periódico \qualisA{} exige:
\begin{itemize}
\item Originalidade comprovada da contribuição;
\item Rigor metodológico na condução da pesquisa;
\item Revisão por pares cega (peer review);
\item Adequação às normas e escopo do periódico;
\item Reproducibilidade dos resultados apresentados.
\end{itemize}

\subsection{Como Agentes de IA Podem Auxiliar na Produção Acadêmica}

O \opencodeeco{} integra 128 agentes especializados que atuam em diferentes
etapas da produção acadêmica \cite{opencode2026ecosystem}:

\begin{itemize}
\item \destaque{Pesquisa bibliográfica}: agentes SEEKER (10 agentes) realizam
      busca sistemática em 10+ fontes acadêmicas (arXiv, OpenAlex, PubMed, CORE,
      Semantic Scholar);
\item \destaque{Escrita científica}: MASWOS v5.0 (49 agentes especializados)
      produz o manuscrito seguindo normas ABNT e template \texttt{abntex2};
\item \destaque{Revisão e correção}: 5 revisores simulados, 4 orientadores PhD,
      6 motores de correção atuam iterativamente;
\item \destaque{Validação estatística}: PhD Auditor aplica testes de Nash,
      Cohen, Bonferroni e análise de sensibilidade;
\item \destaque{Simulação de banca}: Agent Forum debate multiagente simulando
      avaliadores interno, externo e suplente.
\end{itemize}

\subsection{Ética e Boas Práticas: IA como Ferramenta, não Substituta}

\begin{definicao}[Uso Ético de IA na Academia]
A inteligência artificial deve ser utilizada como ferramenta de auxílio à
pesquisa, não como substituta do raciocínio crítico do pesquisador. Todo material
gerado ou assistido por IA deve ser revisado, validado e assumido pelo autor como
de sua inteira responsabilidade \cite{bender2021dangers, gebru2021datasheets}.
\end{definicao}

O \opencodeeco{} incorpora este princípio através de:
\begin{itemize}
\item \destaque{Transparência}: todo texto gerado por agentes é marcado com
      metadados de autoria;
\item \destaque{Auditabilidade}: o raciocínio de cada agente pode ser rastreado
      e verificado;
\item \destaque{Validação humana}: o pesquisador mantém controle editorial final
      sobre todo o conteúdo;
\item \destaque{Detecção de viés}: o motor Critical (15 falácias lógicas)
      analisa o texto em busca de vieses cognitivos e argumentativos.
\end{itemize}

\subsection{Visão Geral do Percurso: Tema $\rightarrow$ Pesquisa $\rightarrow$ Escrita $\rightarrow$ Defesa}

A Figura~\ref{fig:percurso-academico} ilustra as etapas do percurso acadêmico
completo, desde a escolha do tema até a defesa perante a banca.

\begin{figure}[htb]
\centering
\begin{tikzpicture}[node distance=1.8cm, scale=0.4, every node/.style={scale=0.4},
    box/.style={draw, rounded corners, minimum width=2.8cm, minimum height=1.2cm, font=\small, align=center},
    arrow/.style={->, thick},
    label/.style={font=\footnotesize, text=gray!60!black}]
\node[box, fill=blue!10] (tema) {Tema};
\node[box, fill=green!10, right=of tema] (pesq) {Pesquisa};
\node[box, fill=orange!10, right=of pesq] (escrita) {Escrita};
\node[box, fill=red!10, right=of escrita] (correcao) {Correção};
\node[box, fill=purple!10, right=of correcao] (valid) {Validação};
\node[box, fill=cyan!10, right=of valid] (defesa) {Defesa};

\draw[arrow] (tema) -- (pesq);
\draw[arrow] (pesq) -- (escrita);
\draw[arrow] (escrita) -- (correcao);
\draw[arrow] (correcao) -- (valid);
\draw[arrow] (valid) -- (defesa);

\node[label, below=0.3cm of tema] {Nível 0};
\node[label, below=0.3cm of pesq] {Básico};
\node[label, below=0.3cm of escrita] {Intermediário};
\node[label, below=0.3cm of correcao] {Avançado};
\node[label, below=0.3cm of valid] {PhD};
\node[label, below=0.3cm of defesa] {PhD+};
\end{tikzpicture}
\caption{Percurso acadêmico completo: do tema à defesa}
\label{fig:percurso-academico}
\end{figure}

\begin{exercicio}[Nivel 0 -- Reflexivo]
    Pesquise a definição de \qualisA{} para sua área do conhecimento no site da
    CAPES. Liste 5 periódicos A1 da sua área e identifique o escopo editorial
    de cada um.
\end{exercicio}

\begin{exercicio}[Nivel Básico -- Prático]
    Instale o \opencodeeco{} e execute o comando \lstinline{opencode /artigo}
    com o prompt ``produção acadêmica com agentes de IA''. Analise o plano
    gerado pelo MASWOS e identifique as etapas propostas.
\end{exercicio}

% =============================================================================
%  SEÇÃO 8.2: METODOLOGIA PPGTE/UFC: ESTRUTURA DA DISSERTAÇÃO
%  Nível: Avançado (~10 laudas)
% =============================================================================
\section{Metodologia PPGTE/UFC: Estrutura da Dissertação}
\label{sec:metodologia-ppgte-ufc}
\nivelavancado

\paragraph{Os alicerces de uma dissertação sólida} O Programa de Pós-Graduação em Tecnologia Educacional (PPGTE) da Universidade
Federal do Ceará (UFC) estabelece diretrizes específicas para a elaboração de
dissertações, que servem como referência para este capítulo. A estrutura aqui
apresentada é adaptável a outros programas, respeitando-se as particularidades
de cada instituição \cite{abnt2023nbr14724}.

\subsection{Estrutura Padrão da Dissertação}

A dissertação no formato ABNT divide-se em três grandes blocos:

\begin{enumerate}
\item \destaque{Elementos Pré-Textuais} (páginas iniciais sem numeração):
    \begin{itemize}
    \item Capa (obrigatório);
    \item Folha de rosto (obrigatório);
    \item Ficha catalográfica (obrigatório);
    \item Dedicatória (opcional);
    \item Agradecimentos (opcional);
    \item Epígrafe (opcional);
    \item Resumo em português (obrigatório);
    \item Abstract em inglês (obrigatório);
    \item Lista de figuras (opcional);
    \item Lista de tabelas (opcional);
    \item Lista de abreviaturas e siglas (opcional);
    \item Sumário (obrigatório).
    \end{itemize}
\item \destaque{Elementos Textuais} (numeração contínua):
    \begin{itemize}
    \item Introdução (capítulo 1);
    \item Referencial teórico (capítulo 2);
    \item Metodologia (capítulo 3);
    \item Resultados (capítulo 4);
    \item Discussão (capítulo 5);
    \item Conclusão (capítulo 6).
    \end{itemize}
\item \destaque{Elementos Pós-Textuais}:
    \begin{itemize}
    \item Referências bibliográficas (obrigatório);
    \item Apêndices (opcional);
    \item Anexos (opcional);
    \item Glossário (opcional);
    \item Índice remissivo (opcional).
    \end{itemize}
\end{enumerate}

\subsection{Adaptação para Engenharia de Software com Agentes}

Quando a dissertação aborda engenharia de software com agentes inteligentes ---
como é o caso do \opencodeeco{} --- a estrutura tradicional deve ser adaptada
para contemplar:

\begin{itemize}
\item \destaque{Arquitetura do sistema}: descrição detalhada dos componentes,
      suas interações e os padrões de projeto empregados;
\item \destaque{Implementação}: detalhes técnicos da implementação, incluindo
      linguagens, frameworks e ferramentas utilizadas;
\item \destaque{Validação experimental}: experimentos controlados que demonstram
      a eficácia da solução proposta;
\item \destaque{Reprodutibilidade}: disponibilização do código-fonte e dados
      experimentais para verificação independente.
\end{itemize}

\subsection{O Template abntex2 e Personalizações}

O \opencodeeco{} utiliza a classe \codigo{abntex2} \cite{abnt2023nbr14724}
para produção de documentos acadêmicos em conformidade com as normas ABNT.
A classe oferece:

\begin{itemize}
\item Formatação automática de margens (3cm superior/esquerda, 2cm inferior/direita);
\item Espaçamento 1,5 entre linhas configurável;
\item Numeração progressiva de seções;
\item Geração automática de sumário, listas de figuras e tabelas;
\item Suporte a citações no formato autor-data;
\item Compatibilidade com \codigo{biber} para processamento bibliográfico.
\end{itemize}

O template do ecossistema estende a classe base com personalizações específicas:

\begin{lstlisting}[style=bashstyle, caption={Personalizacoes do abntex2 no OpenCode}]
% No preambulo do documento:
\documentclass[12pt,a4paper,oneside,openright]{abntex2}
\usepackage{fontspec}
\usepackage[brazil]{babel}
\usepackage{csquotes}

% Comandos personalizados do ecossistema
\newcommand{\opencode}{\textsc{OpenCode}}
\newcommand{\opencodeeco}{\textsc{OpenCode Ecosystem}}
\newcommand{\qualisA}{\textit{Qualis A1}}
\end{lstlisting}

\subsection{A Dissertação do OpenCode Ecosystem como Estudo de Caso}

A dissertação que documenta o \opencodeeco{} --- intitulada \textit{``OpenCode
Ecosystem v5.4.0: Arquitetura, Implementação e Validação de um Sistema de
Engenharia de Software com Metacognição Funcional e Behavioral Gate
Preventivo''} --- serve como estudo de caso para este capítulo
\cite{opencode2026ecosystem}.

A estrutura da dissertação segue o padrão:

\begin{itemize}
\item \destaque{Capítulo 1}: Fundamentos Matemáticos e Estatísticos ---
      estabelece a base formal para engenharia de software com IA;
\item \destaque{Capítulo 2}: Inteligência Artificial e Arquitetura de Agentes
      Cognitivos --- revisão da literatura sobre sistemas multiagentes;
\item \destaque{Capítulo 3}: OpenCode Ecosystem --- Arquitetura e Engenharia
      de Software com Agentes Inteligentes --- descrição da solução proposta;
\item \destaque{Capítulo 4}: Scanner Pipeline e Metacognição --- metodologia
      de análise e evolução de código;
\item \destaque{Capítulo 5}: Trust Engine e Governança Comportamental ---
      mecanismos de confiança e controle;
\item \destaque{Capítulo 6}: Token Economy e Economia de Agentes --- sistema
      de incentivos econômicos;
\item \destaque{Capítulo 7}: Experimentação e Validação Científica ---
      resultados experimentais com 312 testes passando;
\item \destaque{Capítulo 8}: Dissertação, Produção Científica Qualis A1 e
      Defesa --- validação acadêmica e defesa.
\end{itemize}

A Figura~\ref{fig:estrutura-dissertacao} apresenta a estrutura da dissertação
em formato de diagrama de blocos.

\begin{figure}[htb]
\centering
\begin{tikzpicture}[
    node distance=0.8cm,
    block/.style={draw, rounded corners, minimum width=3cm, minimum height=0.8cm, font=\small, align=center},
    part/.style={draw, rounded corners, minimum width=7cm, minimum height=0.6cm, font=\small\bfseries, align=center, fill=gray!10},
    arrow/.style={->, thick}
]
\node[part, fill=blue!10] (pre) {Parte I -- Fundamentos};
\node[block, below=0.5cm of pre] (c1) {Cap. 1: Matemática e Estatística};
\node[block, below=0.3cm of c1] (c2) {Cap. 2: IA e Agentes Cognitivos};

\node[part, fill=green!10, below=0.8cm of c2] (arq) {Parte II -- Arquitetura};
\node[block, below=0.5cm of arq] (c3) {Cap. 3: OpenCode Ecosystem};
\node[block, below=0.3cm of c3] (c4) {Cap. 4: Scanner e Metacognição};
\node[block, below=0.3cm of c4] (c5) {Cap. 5: Trust Engine};

\node[part, fill=orange!10, below=0.8cm of c5] (exp) {Parte III -- Economia, Experimentação};
\node[block, below=0.5cm of exp] (c6) {Cap. 6: Token Economy};
\node[block, below=0.3cm of c6] (c7) {Cap. 7: Experimentação e Validação};
\node[block, below=0.3cm of c7] (c8) {Cap. 8: Dissertação e Defesa};

\draw[arrow] (c1.south) -- (c2.north);
\draw[arrow] (c2.south) -- (c3.north);
\draw[arrow] (c3.south) -- (c4.north);
\draw[arrow] (c4.south) -- (c5.north);
\draw[arrow] (c5.south) -- (c6.north);
\draw[arrow] (c6.south) -- (c7.north);
\draw[arrow] (c7.south) -- (c8.north);
\end{tikzpicture}
\caption{Estrutura da dissertação do OpenCode Ecosystem}
\label{fig:estrutura-dissertacao}
\end{figure}

\subsection{Elementos Pré-Textuais Detalhados}

A capa da dissertação deve conter: nome da instituição, nome do autor, título,
subtítulo (se houver), número de volumes (se houver), local e ano. A folha de
rosto adiciona a natureza do trabalho (dissertação apresentada ao programa...),
o nome do orientador e a instituição \cite{abnt2023nbr14724}.

A ficha catalográfica, elaborada conforme a Classificação Decimal de Dewey
(CDD), é de responsabilidade da biblioteca da instituição. O \opencodeeco{}
gera automaticamente uma proposta de ficha catalográfica que deve ser validada
pelo bibliotecário responsável.

\subsection{Especificidades do PPGTE/UFC}

O PPGTE/UFC estabelece requisitos adicionais para a dissertação:

\begin{itemize}
\item \destaque{Vínculo com tecnologia educacional}: a pesquisa deve demonstrar
      relação explícita com o campo da tecnologia aplicada à educação;
\item \destaque{Produto educacional}: quando aplicável, a dissertação deve
      incluir um produto educacional (software, sequência didática, guia, etc.);
\item \destaque{Comitê de ética}: pesquisas envolvendo seres humanos devem ser
      submetidas ao Comitê de Ética em Pesquisa (CEP) via Plataforma Brasil;
\item \destaque{Idioma}: a dissertação pode ser redigida em português ou inglês,
      com resumo obrigatório em ambos os idiomas;
\item \destaque{Prazo de entrega}: a versão final deve ser depositada na
      secretaria do programa até 60 dias após a defesa.
\end{itemize}

\begin{exercicio}[Nivel Avançado -- Pesquisa]
    Acesse o site do PPGTE/UFC e identifique as normas complementares para
    dissertações. Compare com as normas gerais da ABNT e liste as diferenças
    específicas do programa.
\end{exercicio}

\begin{exercicio}[Nivel Avançado -- Prático]
    Utilize o template \texttt{abntex2}{} do \opencodeeco{} para gerar um documento
    com todos os elementos pré-textuais preenchidos com dados fictícios.
    Compile e verifique a formatação.
\end{exercicio}

% =============================================================================
%  SEÇÃO 8.3: PROTOCOLO DE ANONIMATO PARA AVALIAÇÃO CEGA
%  Nível: Avançado (~6 laudas)
% =============================================================================
\section{Protocolo de Anonimato para Avaliação Cega}
\label{sec:protocolo-anonimato}
\nivelavancado

\paragraph{A arte invisível de proteger a identidade do autor} A avaliação cega (\textit{blind review}) é um dos pilares do rigor científico
na revisão por pares. No contexto de dissertações e artigos acadêmicos, o
anonimato garante que o mérito do trabalho seja avaliado independentemente da
identidade do autor \cite{popper2005logica, kuhn1998estrutura}.

\subsection{Importância do Anonimato em Avaliação Acadêmica}

A revisão cega serve a múltiplos propósitos:

\begin{enumerate}
\item \destaque{Imparcialidade}: elimina vieses conscientes ou inconscientes
      relacionados à identidade, gênero, raça ou instituição do autor;
\item \destaque{Mérito científico}: a avaliação concentra-se exclusivamente no
      conteúdo, metodologia e resultados do trabalho;
\item \destaque{Credibilidade}: aumenta a confiança da comunidade científica no
      processo de revisão;
\item \destaque{Equidade}: pesquisadores de instituições menos conhecidas têm
      as mesmas oportunidades de publicação.
\end{enumerate}

\subsection{Identificadores Diretos vs. Indiretos}

O protocolo de anonimato classifica os identificadores em duas categorias:

\begin{definicao}[Identificador Direto]
Informação que explicitamente nomeia o autor ou coautores do trabalho.
Exemplos: nome completo, iniciais, ORCID, e-mail institucional, número de
matrícula.
\end{definicao}

\begin{definicao}[Identificador Indireto]
Informação que, combinada com outras fontes, permite inferir a identidade do
autor. Exemplos: nome do orientador, nome do programa de pós-graduação,
agradecimentos a colegas específicos, referência a trabalhos anteriores do
próprio autor, nome da cidade ou instituição.
\end{definicao}

A Tabela~\ref{tab:identificadores} apresenta exemplos de cada tipo e as ações
de anonimato correspondentes.

\begin{table}[htb]
\centering
\caption{Identificadores diretos e indiretos e ações de anonimato}
\label{tab:identificadores}
\small
    \resizebox{\textwidth}{!}{
\begin{tabular}{lp{4.5cm}p{4.5cm}}
\toprule
\textbf{Tipo} & \textbf{Exemplo} & \textbf{Ação} \\
\midrule
Direto & ``Marcelo Claro Laranjeira'' & Substituir por ``Autor'' \\
Direto & \texttt{marcelo@email.com} & Remover ou substituir \\
Direto & ``meu\_lattes.cnpq.br'' & Remover \\
Indireto & ``Prof. Dr. João Silva (orientador)'' & Substituir por ``Orientador'' \\
Indireto & ``Universidade Federal do Ceará'' & Substituir por ``Instituição'' \\
Indireto & ``Agradeço a Maria Souza'' & Remover ou generalizar \\
Indireto & ``como mostramos em (Laranjeira, 2024)'' & Substituir por ``como mostrado em trabalho anterior'' \\
\bottomrule
\end{tabular}
    }
\end{table}

\subsection{Ferramentas de Detecção e Remoção}

O \opencodeeco{} implementa o \destaque{Protocolo de Anonimato} como um
módulo integrado ao pipeline de produção acadêmica. As ferramentas incluem:

\begin{itemize}
\item \destaque{Detector de identificadores diretos}: expressões regulares
      para nomes próprios, e-mails, ORCID, URLs de Lattes e GitHub;
\item \destaque{Detector de identificadores indiretos}: análise contextual
      para identificar instituições, orientadores, agradecimentos;
\item \destaque{Validador de anonimato}: verifica se o documento está
      efetivamente anônimo, reportando riscos residuais;
\item \destaque{Relatório de anonimização}: documento que lista todas as
      substituições realizadas para auditoria posterior.
\end{itemize}

\subsection{Protocolo Implementado no Ecossistema}

O protocolo é executado em 5 etapas:

\begin{enumerate}
\item \destaque{Varredura inicial}: o scanner percorre todo o documento LaTeX
      identificando padrões de identificadores diretos e indiretos;
\item \destaque{Classificação}: cada identificador é classificado como direto
      (prioridade crítica) ou indireto (prioridade alta);
\item \destaque{Substituição automática}: identificadores diretos são
      substituídos automaticamente por termos genéricos;
\item \destaque{Revisão manual assistida}: identificadores indiretos são
      apresentados ao autor para decisão de substituição;
\item \destaque{Validação final}: o validador de anonimato confirma que o
      documento está apto para submissão cega.
\end{enumerate}

\begin{lstlisting}[style=bashstyle, caption={Execucao do protocolo de anonimato}]
# Executa o protocolo de anonimato no documento
opencode /reversa anonymize --input=dissertacao.tex --output=dissertacao_anonima.tex

# Valida o anonimato
opencode /reversa validate-anonymity --input=dissertacao_anonima.tex
\end{lstlisting}

\begin{exercicio}[Nivel Avançado -- Prático]
    Utilize o protocolo de anonimato do \opencodeeco{} em um arquivo LaTeX
    de sua autoria. Execute o validador e identifique quantos identificadores
    foram encontrados e substituídos.
\end{exercicio}

\begin{exercicio}[Nivel Avançado -- Pesquisa]
    Pesquise na literatura (Google Scholar, Scopus) sobre viés em revisão por
    pares. Identifique três estudos que demonstram a eficácia da revisão cega
    na redução de vieses.
\end{exercicio}

% =============================================================================
%  SEÇÃO 8.4: SIMULAÇÃO DE BANCA COM AGENT-FORUM
%  Nível: PhD (~12 laudas)
% =============================================================================
\section{Simulação de Banca com Agent-Forum}
\label{sec:simulacao-banca-agent-forum}
\nivelphd

\paragraph{Simulando o grande dia com múltiplas inteligências} A simulação de banca examinadora é uma das funcionalidades mais avançadas do
\opencodeeco{}, implementada através do módulo Agent Forum (P14-P18) integrado
ao pipeline MiroFish/BettaFish \cite{opencode2026ecosystem, opencode2026spec036}.
Esta seção apresenta a arquitetura, implementação e resultados da simulação.

\subsection{O que é uma Banca Examinadora de Dissertação}

A banca examinadora é um colegiado de professores doutores responsável por
avaliar a dissertação e arguir o candidato. A composição típica inclui:

\begin{itemize}
\item \destaque{Orientador}: presidente da banca, responsável pela mediação;
\item \destaque{Avaliador interno}: professor do mesmo programa de
      pós-graduação;
\item \destaque{Avaliador externo}: professor de outra instituição, garantindo
      independência da avaliação;
\item \destaque{Suplente}: membro que substitui um dos avaliadores em caso de
      impedimento.
\end{itemize}

A arguição segue um ritual acadêmico: o candidato apresenta o trabalho em 20-30
minutos, seguido pela arguição de cada membro (30-60 minutos no total), e
concluído com a deliberação da banca.

\subsection{Agent Forum: Debate Multiagente Simulando Banca}

O Agent Forum do \opencodeeco{} implementa um debate multiagente onde cada
agente assume o papel de um membro da banca \cite{wooldridge2009multiagent}.
O sistema utiliza 212+ tipos de raciocínio distribuídos em 27 categorias para
simular diferentes perspectivas de avaliação.

\begin{figure}[htb]
\centering
\begin{tikzpicture}[
    node distance=2cm,
    agent/.style={draw, circle, minimum size=1.8cm, font=\small\bfseries, align=center},
    forum/.style={draw, rounded corners, minimum width=4cm, minimum height=1cm, font=\small, fill=gray!10},
    arrow/.style={<->, thick, dashed}
]
\node[agent, fill=blue!20] (interno) {Avaliador\\Interno};
\node[agent, fill=green!20, right=of interno] (externo) {Avaliador\\Externo};
\node[agent, fill=red!20, right=of externo] (suplente) {Suplente};
\node[forum, below=2.5cm of externo] (candidato) {Candidato (respostas)};

\draw[arrow] (interno.south) -- (candidato.north west);
\draw[arrow] (externo.south) -- (candidato.north);
\draw[arrow] (suplente.south) -- (candidato.north east);

\node[font=\footnotesize, above=0.2cm of interno] {Persona 1};
\node[font=\footnotesize, above=0.2cm of externo] {Persona 2};
\node[font=\footnotesize, above=0.2cm of suplente] {Persona 3};
\end{tikzpicture}
\caption{Arquitetura do Agent Forum para simulação de banca}
\label{fig:agent-forum-banca}
\end{figure}

\subsection{3 Personas de Banca}

O Agent Forum instancia três personas para a simulação \cite{opencode2026spec036}:

\begin{enumerate}
\item \destaque{Avaliador Interno (Persona 1)}:
    \begin{itemize}
    \item Perfil: pesquisador sênior da mesma área, foco em fundamentação
          teórica e alinhamento com a linha de pesquisa do programa;
    \item Perguntas típicas: ``Qual a contribuição original desta pesquisa?'',
          ``Como este trabalho se relaciona com as pesquisas anteriores do
          grupo?'';
    \item Estilo: construtivo, busca fortalecer a base teórica.
    \end{itemize}
\item \destaque{Avaliador Externo (Persona 2)}:
    \begin{itemize}
    \item Perfil: especialista de outra instituição, foco em metodologia e
          validação experimental;
    \item Perguntas típicas: ``A metodologia adotada é adequada para responder
          à pergunta de pesquisa?'', ``Os resultados são estatisticamente
          significativos?'';
    \item Estilo: crítico, busca identificar fragilidades metodológicas.
    \end{itemize}
\item \destaque{Suplente (Persona 3)}:
    \begin{itemize}
    \item Perfil: pesquisador de área correlata, foco em impacto e
          aplicabilidade;
    \item Perguntas típicas: ``Qual o impacto prático desta pesquisa?'',
          ``Como os resultados podem ser replicados por outros pesquisadores?'';
    \item Estilo: abrangente, busca conexões interdisciplinares.
    \end{itemize}
\end{enumerate}

\subsection{16 Perguntas Simuladas}

O Agent Forum gerou 16 perguntas que cobrem as dimensões essenciais da
avaliação acadêmica. A Tabela~\ref{tab:perguntas-banca} apresenta as perguntas
organizadas por categoria.

\begin{table}[htb]
\centering
    \small
\caption{Perguntas simuladas pelo Agent Forum}
\label{tab:perguntas-banca}
    \resizebox{\textwidth}{!}{
\begin{tabular}{cp{10cm}}
\toprule
\textbf{\#} & \textbf{Pergunta} \\
\midrule
1 & Qual a contribuição original desta pesquisa para a área de engenharia de software? \\
2 & Como você define metacognição funcional e qual a diferença para metacognição humana? \\
3 & A metodologia experimental com 312 testes é suficiente para validar o ecossistema? \\
4 & Como o Trust Engine garante que agentes não ajam de forma maliciosa? \\
5 & Qual o custo computacional do Behavioral Gate e como ele escala? \\
6 & A Token Economy realmente cria incentivos alinhados ou pode ser manipulada? \\
7 & Como você diferencia seu trabalho de outras plataformas de agentes (LangChain, AutoGPT)? \\
8 & Quais as limitações do CORA-Eval como framework de benchmark? \\
9 & Como o protocolo de anonimato foi validado empiricamente? \\
10 & O ciclo evolutivo R1-R23 pode ser replicado por outros pesquisadores? \\
11 & A correção de Bonferroni foi aplicada corretamente considerando as dependências entre testes? \\
12 & Qual a generalidade do ecossistema para domínios além da engenharia de software? \\
13 & Como o Manus Evolve aprende com os ciclos anteriores sem overfitting? \\
14 & O Self-Model N0-N3 é uma forma de consciência artificial? \\
15 & Como a governança cooperativa (Ostrom DP1-DP8) se aplica a agentes de software? \\
16 & Quais as implicações éticas de um sistema que evolve autonomamente seu código? \\
\bottomrule
\end{tabular}
    }
\end{table}

\subsection{Estratégias de Defesa: 6 Estratégias, 8 Configurações}

O Agent Forum implementa 6 estratégias de defesa, cada uma com configurações
ajustáveis \cite{shoham2023multiagent, nash1950equilibrium}:

\begin{enumerate}
\item \destaque{Refutação Direta}: contra-argumentação baseada em evidências
      empíricas do próprio trabalho;
\item \destaque{Redirecionamento Contextual}: reenquadramento da pergunta em
      um contexto mais amplo onde a contribuição é mais clara;
\item \destaque{Admissão com Mitigação}: reconhecimento da limitação seguido
      de proposta de solução ou trabalho futuro;
\item \destaque{Referência Cruzada}: remissão a seções específicas da
      dissertação que tratam do ponto questionado;
\item \destaque{Analogia Estruturada}: uso de analogias com sistemas
      consolidados para explicar conceitos complexos;
\item \destaque{Decomposição Lógica}: quebra da objeção em subproblemas
      respondidos individualmente.
\end{enumerate}

As 8 configurações do Agent Forum permitem ajustar:
\begin{itemize}
\item Nível de agressividade do avaliador (1-5);
\item Profundidade da arguição (superficial a exaustiva);
\item Domínio de conhecimento (estrito ao trabalho ou abrangente);
\item Tolerância a respostas parciais;
\item Estilo de feedback (apenas crítico ou construtivo);
\item Peso relativo de cada dimensão avaliada;
\item Número de iterações do debate;
\item Critério de convergência (score, tempo, ou iterações).
\end{itemize}

\subsection{Nota DAP: 8,07 $\rightarrow$ 9,0 (após Refinamento)}

A \destaque{Nota DAP (Desempenho em Arguição Programática)} é a métrica que
quantifica a qualidade das respostas do candidato durante a simulação. A nota
varia de 0 a 10 e é calculada pela média ponderada de 5 dimensões:

\begin{equation}
DAP = \frac{w_1 \cdot C + w_2 \cdot P + w_3 \cdot F + w_4 \cdot A + w_5 \cdot E}{\sum w_i}
\end{equation}

Onde:
\begin{itemize}
\item $C$ = Correção técnica (peso 3);
\item $P$ = Profundidade da resposta (peso 2);
\item $F$ = Fluência argumentativa (peso 2);
\item $A$ = Adequação à pergunta (peso 2);
\item $E$ = Evidências apresentadas (peso 1).
\end{itemize}

A Tabela~\ref{tab:evolucao-dap} mostra a evolução da nota DAP ao longo dos
ciclos de refinamento.

\begin{table}[htb]
\centering
    \small
\caption{Evolução da Nota DAP durante o refinamento}
\label{tab:evolucao-dap}
    \resizebox{\textwidth}{!}{
\begin{tabular}{ccccccc}
\toprule
\textbf{Ciclo} & \textbf{C} & \textbf{P} & \textbf{F} & \textbf{A} & \textbf{E} & \textbf{DAP} \\
\midrule
1 & 7,5 & 7,0 & 8,0 & 8,5 & 9,0 & 7,80 \\
2 & 8,0 & 7,5 & 8,0 & 8,5 & 9,0 & 8,07 \\
3 & 8,5 & 8,0 & 8,5 & 9,0 & 9,5 & 8,55 \\
4 & 9,0 & 8,5 & 9,0 & 9,0 & 9,5 & 8,90 \\
5 & 9,0 & 9,0 & 9,0 & 9,0 & 10,0 & 9,00 \\
\bottomrule
\end{tabular}
    }
\end{table}

A evolução de 8,07 para 9,0 demonstra a eficácia do processo iterativo de
refinamento: cada ciclo identifica fragilidades nas respostas e propõe
melhorias específicas.

\subsection{212+ Tipos de Raciocínio Aplicados à Defesa}

O ecossistema dispõe de 212+ tipos de raciocínio distribuídos em 27 categorias
\cite{opencode2026ecosystem}, dos quais os seguintes são mais relevantes para
a defesa acadêmica:

\begin{itemize}
\item \destaque{Raciocínio Lógico (5 subtipos)}: dedução, indução, abdução,
      silogismo, contraposição;
\item \destaque{Raciocínio Dialético (5 subtipos)}: tese-antítese-síntese,
      refutação, concessão, redirecionamento, reconciliação;
\item \destaque{Teoria dos Jogos (10 subtipos)}: equilíbrio de Nash, estratégia
      dominante, min-max, leilão, barganha;
\item \destaque{Raciocínio Probabilístico (8 subtipos)}: Bayesiano,
      frequentista, Monte Carlo, bootstrap, sensibilidade;
\item \destaque{Argumentação Estruturada (12 subtipos)}: Toulmin,
      argumentação pragmática, ética, analógica, causal.
\end{itemize}

A Figura~\ref{fig:raciocinios-defesa} apresenta a hierarquia dos tipos de
raciocínio aplicados à defesa acadêmica, organizados por categoria e nível
de complexidade.

\begin{figure}[htb]
\centering
\begin{tikzpicture}[
    level distance=1.2cm,
    sibling distance=3cm,
    every node/.style={draw, rounded corners, font=\tiny, align=center},
    edge from parent/.style={draw, ->, thick}
]
\node[fill=red!20, minimum width=2cm, minimum height=0.6cm] {212+ Raciocínios}
    child {node[fill=orange!20, minimum width=1.8cm, minimum height=0.5cm] {Lógico (5)}
        child {node[fill=yellow!20, minimum width=1.5cm] {Dedução}}
        child {node[fill=yellow!20, minimum width=1.5cm] {Indução}}
        child {node[fill=yellow!20, minimum width=1.5cm] {Abdução}}
    }
    child {node[fill=green!20, minimum width=1.8cm, minimum height=0.5cm] {Dialético (5)}
        child {node[fill=yellow!20, minimum width=1.5cm] {Tese}}
        child {node[fill=yellow!20, minimum width=1.5cm] {Antítese}}
        child {node[fill=yellow!20, minimum width=1.5cm] {Síntese}}
    }
    child {node[fill=blue!20, minimum width=1.8cm, minimum height=0.5cm] {Jogos (10)}
        child {node[fill=yellow!20, minimum width=1.5cm] {Nash}}
        child {node[fill=yellow!20, minimum width=1.5cm] {Min-Max}}
        child {node[fill=yellow!20, minimum width=1.5cm] {Barganha}}
    }
    child {node[fill=purple!20, minimum width=1.8cm, minimum height=0.5cm] {Probab. (8)}
        child {node[fill=yellow!20, minimum width=1.5cm] {Bayes}}
        child {node[fill=yellow!20, minimum width=1.5cm] {Monte Carlo}}
        child {node[fill=yellow!20, minimum width=1.5cm] {Sensibilidade}}
    };
\end{tikzpicture}
\caption{Hierarquia dos 212+ tipos de raciocínio aplicados à defesa acadêmica}
\label{fig:raciocinios-defesa}
\end{figure}

Além das categorias listadas, o ecossistema aplica raciocínios avançados como:

\begin{itemize}
\item \destaque{Raciocínio Causal (Granger, Pearl)}: identificação de relações
      causais entre variáveis do experimento, fundamental para responder a
      perguntas sobre validade interna;
\item \destaque{Raciocínio Metacognitivo (Flavell, Schraw)}: reflexão sobre o
      próprio processo de raciocínio, permitindo ao agente identificar quando
      uma resposta é insuficiente;
\item \destaque{Raciocínio Cooperativo (Ostrom)}: modelagem de interações
      cooperativas entre múltiplos agentes, aplicável a perguntas sobre
      governança do ecossistema;
\item \destaque{Raciocínio Ético (Rawls, Kant)}: avaliação de implicações
      éticas das decisões de projeto, relevante para perguntas sobre
      responsabilidade e viés algorítmico.
\end{itemize}

A combinação destes raciocínios em cascata --- onde um raciocínio lógico
fornece a estrutura, um raciocínio dialético explora contradições, e um
raciocínio probabilístico quantifica a incerteza --- produz respostas
robustas e multifacetadas, características de uma defesa de nível PhD.

\begin{exercicio}[Nivel PhD -- Simulação]
    Execute o Agent Forum do \opencodeeco{} com o comando
    \lstinline{opencode /auto agent-forum --mode=banca}. Analise as perguntas
    geradas e classifique cada uma segundo as 6 estratégias de defesa.
\end{exercicio}

\begin{exercicio}[Nivel PhD -- Reflexivo]
    Escolha 3 perguntas da Tabela~\ref{tab:perguntas-banca} e elabore
    respostas completas utilizando cada uma das 6 estratégias de defesa.
    Compare a eficácia de cada estratégia.
\end{exercicio}

% =============================================================================
%  SEÇÃO 8.5: PhD AUDITOR: NASH, COHEN, BONFERRONI, QUALIS
%  Nível: PhD (~10 laudas)
% =============================================================================
\section{PhD Auditor: Nash, Cohen, Bonferroni, Qualis}
\label{sec:phd-auditor}
\nivelphd

\paragraph{O rigor estatístico como guardião da verdade científica} O \destaque{PhD Auditor} é o módulo de validação científica de mais alto nível
do \opencodeeco{}, responsável por aplicar rigor estatístico e formal à
avaliação da dissertação \cite{opencode2026ecosystem, opencode2026spec036}.
Ele integra quatro motores especializados que operam de forma coordenada.

\subsection{Nash Solver: Equilíbrio de Nash em Revisão por Pares}

O Nash Solver modela a revisão por pares como um jogo cooperativo onde
revisores e autores buscam maximizar a qualidade científica do trabalho
\cite{nash1950equilibrium, vonneumann1944theory}.

\begin{definicao}[Jogo da Revisão por Pares]
Seja $N = \{1, \dots, n\}$ o conjunto de revisores e $A = \{a_1, \dots, a_m\}$
o conjunto de ações possíveis (aprovar, aprovar com correções, rejeitar,
solicitar revisão maior). Cada revisor $i$ possui uma função utilidade
$U_i(a_1, \dots, a_n)$ que depende das ações de todos os revisores.
O equilíbrio de Nash ocorre quando nenhum revisor pode aumentar sua utilidade
unilateralmente \cite{nash1950equilibrium}:
\begin{equation}
U_i(a_i^*, a_{-i}^*) \geq U_i(a_i, a_{-i}^*), \quad \forall i \in N, \forall a_i \in A
\end{equation}
\end{definicao}

O Nash Solver implementa um algoritmo iterativo que:

\begin{enumerate}
\item Inicializa as ações de todos os revisores (aprovação condicional);
\item Para cada revisor, calcula a melhor resposta dadas as ações dos demais;
\item Atualiza as ações até convergência (nenhum revisor deseja mudar);
\item Retorna o equilíbrio encontrado e a pontuação de qualidade associada.
\end{enumerate}

A convergência para o equilíbrio de Nash na simulação de banca do \opencodeeco{}
ocorreu em média após 4,7 iterações (desvio padrão de 1,2), indicando que o
sistema atinge estabilidade rapidamente.

\subsection{Statistical Rigor (Cohen): Tamanho de Efeito e Poder Estatístico}

O Statistical Rigor module aplica os princípios de Jacob Cohen para garantir
que as análises estatísticas da dissertação atendam aos padrões de rigor
científico \cite{cohen2018estatistica, wasserman2004all}.

\begin{definicao}[Tamanho de Efeito (d de Cohen)]
O \destaque{d de Cohen} quantifica a magnitude da diferença entre dois grupos,
independentemente do tamanho da amostra:
\begin{equation}
d = \frac{\bar{x}_1 - \bar{x}_2}{s_p}, \quad s_p = \sqrt{\frac{(n_1-1)s_1^2 + (n_2-1)s_2^2}{n_1 + n_2 - 2}}
\end{equation}
Onde $\bar{x}_i$ são as médias, $s_i$ os desvios padrão, e $n_i$ os tamanhos
das amostras. Valores de referência: $d=0,2$ (pequeno), $d=0,5$ (médio),
$d=0,8$ (grande) \cite{cohen2018estatistica}.
\end{definicao}

O calculador de poder estatístico determina o tamanho amostral necessário para
detectar um efeito de magnitude esperada com determinada confiança:

\begin{equation}
n = \frac{2(Z_{\alpha/2} + Z_\beta)^2 \sigma^2}{\delta^2}
\end{equation}

Onde $\alpha$ é o nível de significância (tipicamente 0,05), $\beta$ é a
probabilidade de erro tipo II (tipicamente 0,20, poder = 0,80), $\sigma$ é o
desvio padrão esperado e $\delta$ é a diferença mínima detectável.

Para a validação do \opencodeeco{}, o Statistical Rigor calculou:

\begin{itemize}
\item \destaque{d de Cohen} para a comparação entre ciclos evolutivos: $d = 0,73$ (efeito médio-alto);
\item \destaque{Poder estatístico} para a bateria de 312 testes: $1-\beta = 0,94$ (excelente);
\item \destaque{Tamanho amostral mínimo} recomendado: $n \geq 30$ por grupo
      (atendido com folga pela suíte de 312 testes).
\end{itemize}

\subsection{Bonferroni Correction: Múltiplas Comparações}

Quando múltiplas hipóteses são testadas simultaneamente, a probabilidade de
falsos positivos (erro tipo I) aumenta. A correção de Bonferroni ajusta o
nível de significância para controlar a \destaque{Família de Erros}
(Family-Wise Error Rate, FWER) \cite{bonferroni1936teoria}.

\begin{definicao}[Correção de Bonferroni]
Sejam $m$ hipóteses sendo testadas simultaneamente. A correção de Bonferroni
ajusta o nível de significância individual para:
\begin{equation}
\alpha_{\text{individual}} = \frac{\alpha_{\text{global}}}{m}
\end{equation}
Uma hipótese nula é rejeitada apenas se seu p-valor for menor que
$\alpha_{\text{individual}}$ \cite{bonferroni1936teoria}.
\end{definicao}

No contexto do \opencodeeco{}, a correção de Bonferroni foi aplicada nas
seguintes análises:

\begin{itemize}
\item \destaque{Comparação entre ciclos evolutivos}: $m = 22$ comparações
      (R1 vs R2, R1 vs R3, \dots, R22 vs R23), com $\alpha_{\text{global}} = 0,05$,
      resultando em $\alpha_{\text{individual}} = 0,0023$;
\item \destaque{Dimensões do CORA-Eval}: $m = 10$ dimensões, com
      $\alpha_{\text{individual}} = 0,005$;
\item \destaque{Scanners do pipeline}: $m = 5$ scanners, com
      $\alpha_{\text{individual}} = 0,01$.
\end{itemize}

\subsection{Qualis A1 Auditor: Verificação Automática dos Critérios}

O Qualis A1 Auditor implementa a verificação automática dos critérios
necessários para classificação de um periódico como \qualisA{}
\cite{laranjeira2026opencode, laranjeira2026gartner}.

Os critérios verificados incluem:

\begin{enumerate}
\item \destaque{Indexação}: o periódico está indexado nas bases Web of Science,
      Scopus, ou SciELO?
\item \destaque{Fator de impacto}: o JCR (Journal Citation Reports) está acima
      do percentil 75 da área?
\item \destaque{Corpo editorial}: o periódico possui corpo editorial
      internacional e diversificado?
\item \destaque{Revisão por pares}: o processo de revisão é cego e
      sistematizado?
\item \destaque{Aceitação}: a taxa de aceitação é inferior a 30\%?
\item \destaque{Tempo de publicação}: o tempo médio entre submissão e
      publicação é adequado?
\item \destaque{Qualidade dos artigos}: os artigos publicados demonstram rigor
      metodológico e originalidade?
\end{enumerate}

Para cada critério, o auditor atribui uma pontuação de 0 a 10 e calcula a
média ponderada. A pontuação mínima para classificação \qualisA{} é 8,0.

\subsection{Sensitivity Analyzer: Análise de Sensibilidade dos Resultados}

O Sensitivity Analyzer realiza análise de sensibilidade para determinar a
robustez dos resultados apresentados \cite{box1976science}. As técnicas
implementadas incluem:

\begin{itemize}
\item \destaque{Análise one-at-a-time (OAT)}: variação de um parâmetro por vez
      para medir seu impacto no resultado;
\item \destaque{Simulação de Monte Carlo}: amostragem aleatória dos parâmetros
      dentro de seus intervalos de confiança;
\item \destaque{Análise de cenários}: definição de cenários otimista, esperado
      e pessimista para cada resultado;
\item \destaque{Índice de sensibilidade de Sobol}: decomposição da variância
      total em contribuições de cada parâmetro.
\end{itemize}

\subsection{IMRAD Formatter: Formatação Introdução-Métodos-Resultados-Discussão}

O IMRAD Formatter garante que a dissertação siga a estrutura Introdução-
Métodos-Resultados-Discussão (IMRAD), padrão internacional para artigos
científicos \cite{sommerville2011engenharia}.

A formatação IMRAD é aplicada automaticamente, assegurando:

\begin{itemize}
\item \destaque{Introdução}: contextualização, problema de pesquisa,
      objetivos e justificativa;
\item \destaque{Métodos}: descrição detalhada dos procedimentos,
      instrumentos e técnicas de análise;
\item \destaque{Resultados}: apresentação objetiva dos dados obtidos, sem
      interpretação;
\item \destaque{Discussão}: interpretação dos resultados à luz da literatura,
      limitações e implicações.
\end{itemize}

\subsection{Como o PhD Auditor Valida a Dissertação}

O fluxo de validação integrado do PhD Auditor segue as etapas:

\begin{enumerate}
\item \destaque{Pré-validação}: verificação de formatação, estrutura e
      conformidade com normas ABNT;
\item \destaque{Validação estatística}: aplicação dos motores Nash, Cohen e
      Bonferroni;
\item \destaque{Validação Qualis}: verificação dos critérios de classificação
      \qualisA{};
\item \destaque{Análise de sensibilidade}: teste de robustez dos resultados;
\item \destaque{Formatação IMRAD}: reestruturação conforme padrão
      internacional;
\item \destaque{Relatório de auditoria}: documento consolidado com hall da
      validação (passou/não passou) para cada critério.
\end{enumerate}

A Figura~\ref{fig:fluxo-phd-auditor} ilustra o fluxo completo.

\begin{figure}[htb]
\centering
\begin{tikzpicture}[
    node distance=1.5cm,
    box/.style={draw, rounded corners, minimum width=3.2cm, minimum height=1cm, font=\small, align=center},
    arrow/.style={->, thick}
]
\node[box, fill=blue!10] (pre) {Pré-validação};
\node[box, fill=green!10, below=of pre] (estat) {Validação\\Estatística};
\node[box, fill=orange!10, below=of estat] (qualis) {Validação\\Qualis};
\node[box, fill=red!10, below=of qualis] (sens) {Análise de\\Sensibilidade};
\node[box, fill=purple!10, below=of sens] (imrad) {Formatação\\IMRAD};
\node[box, fill=cyan!10, below=of imrad] (relat) {Relatório de\\Auditoria};

\draw[arrow] (pre) -- (estat);
\draw[arrow] (estat) -- (qualis);
\draw[arrow] (qualis) -- (sens);
\draw[arrow] (sens) -- (imrad);
\draw[arrow] (imrad) -- (relat);
\end{tikzpicture}
\caption{Fluxo de validação do PhD Auditor}
\label{fig:fluxo-phd-auditor}
\end{figure}

\begin{exercicio}[Nivel PhD -- Pesquisa]
    Pesquise o fator de impacto JCR da sua área de pesquisa. Identifique 3
    periódicos \qualisA{} e analise se eles atendem aos 7 critérios do Qualis
    A1 Auditor.
\end{exercicio}

\begin{exercicio}[Nivel PhD -- Prático]
    Execute o PhD Auditor do \opencodeeco{} em um artigo de sua autoria.
    Analise o relatório gerado e identifique quais critérios foram atendidos
    e quais precisam de melhoria.
\end{exercicio}

% =============================================================================
%  SEÇÃO 8.6: AUTO_SCORE_QUALIS: SISTEMA AUTOMÁTICO DE PONTUAÇÃO
%  Nível: Avançado (~8 laudas)
% =============================================================================
\section{AUTO\_SCORE\_QUALIS: Sistema Automático de Pontuação}
\label{sec:auto-score-qualis}
\nivelavancado

\paragraph{Dez critérios que separam o excepcional do mediano} O \destaque{AUTO\_SCORE\_QUALIS} é o sistema automático de pontuação que
avalia a qualidade da dissertação segundo 10 critérios objetivos, calibrados
por pesos de revisores simulados \cite{opencode2026ecosystem}. O sistema opera
iterativamente: escrever $\rightarrow$ avaliar $\rightarrow$ corrigir
$\rightarrow$ reavaliar, até atingir a pontuação mínima de 95 pontos.

\subsection{10 Critérios de Avaliação}

A Tabela~\ref{tab:criterios-qualis} apresenta os 10 critérios de avaliação
com seus pesos e descrições.

\begin{table}[htb]
\centering
    \small
\caption{Critérios de avaliação do AUTO\_SCORE\_QUALIS}
\label{tab:criterios-qualis}
    \resizebox{\textwidth}{!}{
\begin{tabular}{clc}
\toprule
\textbf{Critério} & \textbf{Descrição} & \textbf{Peso} \\
\midrule
1 & Originalidade da contribuição & 15\% \\
2 & Relevância científica e social & 12\% \\
3 & Qualidade da fundamentação teórica & 12\% \\
4 & Adequação metodológica & 15\% \\
5 & Consistência dos resultados & 12\% \\
6 & Qualidade da discussão & 10\% \\
7 & Clareza e coesão textual & 8\% \\
8 & Conformidade com normas ABNT & 6\% \\
9 & Profundidade das referências & 5\% \\
10 & Reprodutibilidade dos experimentos & 5\% \\
\midrule
& \textbf{Total} & \textbf{100\%} \\
\bottomrule
\end{tabular}
    }
\end{table}

Cada critério é avaliado em uma escala de 0 a 10, e a pontuação final é
calculada como a média ponderada:

\begin{equation}
\text{Score} = \frac{\sum_{i=1}^{10} w_i \cdot s_i}{\sum_{i=1}^{10} w_i} \times 10
\end{equation}

Onde $w_i$ é o peso do critério $i$ e $s_i$ é a pontuação atribuída pelo
revisor.

\subsection{Pesos de Revisores: Calibração Automática}

O sistema utiliza múltiplos revisores simulados, cada um com um perfil de
avaliação distinto. Os pesos de cada revisor são calibrados automaticamente
através de um processo de validação cruzada \cite{hastie2009elements}:

\begin{enumerate}
\item \destaque{Inicialização}: cada revisor recebe pesos iguais para todos
      os critérios;
\item \destaque{Calibração}: um conjunto de dissertações de referência
      (previamente avaliadas por humanos) é utilizado para ajustar os pesos;
\item \destaque{Validação}: a calibração é validada contra um conjunto de
      teste separado;
\item \destaque{Aplicação}: os pesos calibrados são utilizados para avaliar
      novas dissertações.
\end{enumerate}

\subsection{Processo Iterativo: Escrever $\rightarrow$ Avaliar $\rightarrow$ Corrigir $\rightarrow$ Reavaliar}

O ciclo iterativo segue o algoritmo:

\begin{lstlisting}[language=Python, caption={Algoritmo do ciclo iterativo de correcao}]
def ciclo_iterativo(manuscrito, score_minimo=95, max_iteracoes=50):
    score_atual = 0
    historico = []
    for iteracao in range(max_iteracoes):
        # 1. Avaliar manuscrito atual
        score_atual = auto_score_qualis(manuscrito)
        historico.append(score_atual)
        
        # 2. Verificar criterio de parada
        if score_atual >= score_minimo:
            print(f"Score {score_atual} >= {score_minimo}. Convergiu!")
            break
        
        # 3. Identificar criterios abaixo do limiar
        gaps = identificar_gaps(manuscrito, score_atual)
        
        # 4. Corrigir gaps
        for gap in gaps:
            manuscrito = aplicar_correcao(manuscrito, gap)
        
        print(f"Iteracao {iteracao}: score {score_atual:.2f}")
    
    return manuscrito, score_atual, historico
\end{lstlisting}

\subsection{Pontuação: 74 $\rightarrow$ 95 (Evolução Através de Ciclos)}

A Tabela~\ref{tab:evolucao-score} documenta a evolução da pontuação ao longo
dos ciclos de correção.

\begin{table}[htb]
\centering
    \small
\caption{Evolução da pontuação AUTO\_SCORE\_QUALIS}
\label{tab:evolucao-score}
    \resizebox{\textwidth}{!}{
\begin{tabular}{cccc}
\toprule
\textbf{Ciclo} & \textbf{Score} & \textbf{Principais gaps} & \textbf{Ações corretivas} \\
\midrule
1 & 74 & Metodologia, resultados & Revisão da metodologia experimental \\
2 & 78 & Discussão, clareza & Reestruturação da discussão \\
3 & 82 & Referências, ABNT & Expansão e correção das referências \\
4 & 85 & Originalidade, relevância & Fortalecimento da contribuição \\
5 & 88 & Reprodutibilidade & Adição de scripts e dados \\
6 & 90 & Profundidade teórica & Expansão do referencial teórico \\
7 & 92 & Coesão textual & Revisão de coesão e fluência \\
8 & 94 & Consistência resultados & Ajustes finos nos resultados \\
9 & 95 & -- & Convergência alcançada \\
\bottomrule
\end{tabular}
    }
\end{table}

A evolução de 74 para 95 pontos (aumento de 28\%) demonstra a eficácia do
processo iterativo. Cada ciclo identifica e corrige gaps específicos,
elevando progressivamente a qualidade geral do manuscrito.

\begin{figure}[htb]
\centering
\begin{tikzpicture}
\begin{axis}[
    width=12cm, height=6cm,
    xlabel={Ciclo de correção},
    ylabel={Score Qualis},
    xmin=1, xmax=9,
    ymin=70, ymax=100,
    xtick={1,2,3,4,5,6,7,8,9},
    ytick={70,75,80,85,90,95,100},
    grid=both,
    minor grid style={gray!30},
    major grid style={gray!50},
    legend style={at={(0.98,0.02)}, anchor=south east}
]
\addplot[color=blue,mark=*,thick] coordinates {
    (1,74) (2,78) (3,82) (4,85) (5,88) (6,90) (7,92) (8,94) (9,95)
};
\addplot[color=red,dashed,thick] coordinates {(1,95) (9,95)};
\legend{Score, Limiar 95};
\end{axis}
\end{tikzpicture}
\caption{Evolução do score Qualis ao longo dos ciclos de correção}
\label{fig:evolucao-score-qualis}
\end{figure}

\subsection{Arquivo: auto\_score\_qualis.py}

O sistema está implementado no arquivo \codigo{auto\_score\_qualis.py} no
diretório \codigo{criador-artigo/} do ecossistema. A estrutura principal
inclui:

\begin{itemize}
\item \destaque{Classe \texttt{QualisScorer}}: implementa os 10 critérios
      de avaliação e o cálculo da pontuação ponderada;
\item \destaque{Classe \texttt{ReviewerCalibrator}}: realiza a calibração
      automática dos pesos dos revisores;
\item \destaque{Classe \texttt{IterativeCorrector}}: gerencia o ciclo
      iterativo de correção;
\item \destaque{Função \texttt{gap\_analyzer}}: identifica os critérios com
      pontuação abaixo do limiar e sugere ações corretivas;
\item \destaque{Função \texttt{generate\_report}}: produz o relatório
      detalhado de avaliação.
\end{itemize}

\begin{exercicio}[Nivel Avançado -- Prático]
    Execute o AUTO\_SCORE\_QUALIS em um capítulo do livro que você está
    escrevendo. Analise o relatório gerado e identifique os 3 critérios
    com menor pontuação.
\end{exercicio}

\begin{exercicio}[Nivel Avançado -- Programação]
    Modifique o arquivo \codigo{auto\_score\_qualis.py} para adicionar um
    novo critério de avaliação (ex: ``Inovação tecnológica''). Recalibre
    os pesos dos revisores e execute o sistema novamente.
\end{exercicio}

% =============================================================================
%  SEÇÃO 8.7: ITERATIVE CORRECTION LOOP: CICLO DE REFINAMENTO
%  Nível: Avançado (~8 laudas)
% =============================================================================
\section{Iterative Correction Loop: Ciclo de Refinamento}
\label{sec:iterative-correction-loop}
\nivelavancado

\paragraph{O ciclo virtuoso do aperfeiçoamento contínuo} O \destaque{Iterative Correction Loop} é o motor de refinamento contínuo do
\opencodeeco{}, que combina 5 revisores simulados, 4 orientadores PhD e
6 motores de correção para elevar a qualidade do manuscrito
\cite{opencode2026ecosystem, opencode2026maswos}.

\subsection{5 Revisores Simulados}

Cada revisor simulado possui um perfil de avaliação distinto:

\begin{enumerate}
\item \destaque{Revisor Metodológico}: foco em adequação metodológica,
      validade interna e externa, reprodutibilidade;
\item \destaque{Revisor Teórico}: foco em fundamentação teórica, profundidade
      das referências, alinhamento com estado da arte;
\item \destaque{Revisor Estatístico}: foco em análise estatística, testes de
      significância, tamanho de efeito, poder estatístico;
\item \destaque{Revisor Textual}: foco em clareza, coesão, gramática, estilo
      acadêmico, conformidade ABNT;
\item \destaque{Revisor Geral}: visão holística da contribuição, originalidade,
      relevância e impacto.
\end{enumerate}

Cada revisor gera um parecer estruturado com:

\begin{itemize}
\item Pontuação para cada critério (0-10);
\item Justificativa da pontuação;
\item Sugestões de melhoria específicas;
\item Exemplos de trechos que precisam de correção;
\item Prioridade da correção (crítica, alta, média, baixa).
\end{itemize}

\subsection{4 Orientadores/Consultores (PhD)}

Após a avaliação pelos 5 revisores, 4 orientadores/consultores PhD analisam
os pareceres e produzem recomendações consolidadas:

\begin{enumerate}
\item \destaque{Orientador 1 -- Metodologia}: avalia as críticas metodológicas
      e sugere ajustes no design experimental;
\item \destaque{Orientador 2 -- Argumentação}: avalia a solidez da
      argumentação e sugere refinamentos na linha de raciocínio;
\item \destaque{Orientador 3 -- Contribuição}: avalia a originalidade e
      relevância da contribuição, sugerindo fortalecimento;
\item \destaque{Orientador 4 -- Coerência}: avalia a coerência geral do
      manuscrito, identificando contradições ou lacunas.
\end{enumerate}

\subsection{6 Motores de Correção}

Seis motores de correção atuam sobre o manuscrito, cada um especializado em
um aspecto da qualidade textual \cite{opencode2026scanner}:

\begin{enumerate}
\item \destaque{Gramática (Agente 38)}: correção ortográfica, concordância
      verbal e nominal, regência, pontuação;
\item \destaque{Estilo (Agente 39)}: adequação ao estilo acadêmico formal,
      eliminação de coloquialismos, padronização terminológica;
\item \destaque{ABNT (Agente 40)}: verificação de citações, referências,
      formatação de elementos pré e pós-textuais;
\item \destaque{CJK (Agente 41)}: detecção e remoção de caracteres CJK
      (Chinês, Japonês, Coreano) do texto em português;
\item \destaque{Clareza (Agente 42)}: simplificação de parágrafos complexos,
      melhoria da legibilidade, índice Flesch;
\item \destaque{Coesão (Agente 43)}: conectivos, transições entre parágrafos,
      estrutura argumentativa.
\end{enumerate}

\subsection{Correção Textual Qualis (Agente 44)}

O \destaque{Agente 44} (Correção Textual Qualis) é um meta-motor que coordena
os 6 motores de correção e aplica as correções de forma integrada. Ele opera
em três níveis:

\begin{enumerate}
\item \destaque{Nível 1 -- Correção Superficial}: ortografia, gramática,
      formatação ABNT (motores 38, 40);
\item \destaque{Nível 2 -- Correção Estrutural}: estilo, clareza, coesão
      (motores 39, 42, 43);
\item \destaque{Nível 3 -- Correção Profunda}: argumentação, consistência
      interna, alinhamento com objetivos (todos os motores).
\end{enumerate}

\subsection{Refinamento de Argumentação (Agente 45)}

O \destaque{Agente 45} (Refinamento de Argumentação) é especializado em
fortalecer a cadeia argumentativa do manuscrito. Ele utiliza os 212+ tipos
de raciocínio do ecossistema para:

\begin{itemize}
\item Identificar lacunas na argumentação;
\item Sugerir contra-argumentos e refutações;
\item Fortalecer a conexão entre evidências e conclusões;
\item Garantir que cada claim seja suportada por evidência adequada.
\end{itemize}

\subsection{Execução Iterativa até Score $\geq$ 95}

O ciclo completo segue o fluxo:

\begin{enumerate}
\item Manuscrito é submetido aos 5 revisores;
\item Pareceres são consolidados pelos 4 orientadores;
\item Gaps identificados são corrigidos pelos 6 motores;
\item Agente 44 coordena e integra as correções;
\item Agente 45 refina a argumentação;
\item AUTO\_SCORE\_QUALIS reavalia o manuscrito;
\item Se score $\geq$ 95, o ciclo converge; caso contrário, retorna ao passo 1.
\end{enumerate}

A Figura~\ref{fig:ciclo-correcao} ilustra o fluxo completo.

\begin{figure}[htb]
\centering
\resizebox{\textwidth}{!}{%
\begin{tikzpicture}[
    node distance=1.2cm,
    box/.style={draw, rounded corners, minimum width=2.5cm, minimum height=0.8cm, font=\tiny, align=center},
    arrow/.style={->, thick}
]
\node[box, fill=blue!10] (revisores) {5 Revisores\\Simulados};
\node[box, fill=green!10, right=of revisores] (orient) {4 Orientadores\\PhD};
\node[box, fill=orange!10, right=of orient] (motores) {6 Motores\\Correção};
\node[box, fill=red!10, right=of motores] (qualis) {Agente 44\\Coordenação};
\node[box, fill=purple!10, right=of qualis] (arg) {Agente 45\\Argumentação};
\node[box, fill=cyan!10, below=1.5cm of arg] (score) {AUTO\_SCORE\\QUALIS};
\node[box, fill=gray!10, below=1.5cm of score] (fim) {Score $\geq$ 95?\\FIM};

\draw[arrow] (revisores) -- (orient);
\draw[arrow] (orient) -- (motores);
\draw[arrow] (motores) -- (qualis);
\draw[arrow] (qualis) -- (arg);
\draw[arrow] (arg) -- (score);
\draw[arrow] (score) -- (fim);
\draw[arrow] (fim.west) -- ++(-0.5,0) |- (revisores.south);
\draw[arrow] (fim.south) -- ++(0,-0.5);
\end{tikzpicture}%
}
\caption{Ciclo iterativo de correção}
\label{fig:ciclo-correcao}
\end{figure}

\begin{exercicio}[Nivel Avançado -- Simulação]
    Execute o Iterative Correction Loop em um texto de sua autoria (mínimo
    5 páginas). Documente quantas iterações foram necessárias para atingir
    score $\geq$ 95 e quais os principais gaps corrigidos.
\end{exercicio}

\begin{exercicio}[Nivel Avançado -- Programação]
    Implemente um novo motor de correção (Agente 46) especializado em
    detecção de viés de gênero na linguagem acadêmica. Integre-o ao
    Iterative Correction Loop e teste em um corpus de artigos científicos.
\end{exercicio}

% =============================================================================
%  SEÇÃO 8.8: PRODUÇÃO DE ARTIGOS QUALIS A1
%  Nível: PhD (~10 laudas)
% =============================================================================
\section{Produção de Artigos Qualis A1}
\label{sec:producao-artigos-qualis}
\nivelphd

\paragraph{Transformando pesquisa em contribuição reconhecida} A produção de artigos para periódicos \qualisA{} é o coroamento da pesquisa
acadêmica. Esta seção apresenta o pipeline completo de produção de artigos
de alto impacto utilizando o \opencodeeco{} \cite{laranjeira2026gartner,
laranjeira2026cora, laranjeira2026opencode}.

\subsection{Mapeamento Sistemático: Gartner Hype Cycle 2026 vs OpenCode}

O \destaque{Gartner Hype Cycle for Emerging Technologies, 2026} \cite{gartner2026hype}
mapeou 25 tecnologias emergentes, das quais o \opencodeeco{} possui aderência
em 32\% (alta), 20\% (média) e 48\% (baixa). O mapeamento sistemático,
documentado no artigo \textit{``Mapeamento Sistemático do Gartner Hype Cycle
2026 vs. OpenCode Ecosystem''} \cite{laranjeira2026gartner}, utilizou 36
referências e identificou 3 gaps estratégicos:

\begin{enumerate}
\item \destaque{Gap 1 -- Federated API Governance}: governança de APIs
      federadas (SPEC-019, 8 CTs);
\item \destaque{Gap 2 -- Data Streaming Enterprise}: streaming de dados
      empresariais (SPEC-020, 10 CTs);
\item \destaque{Gap 3 -- Low-Code Agent Platform}: plataforma low-code para
      agentes (SPEC-021, 6 CTs).
\end{enumerate}

Cada gap foi endereçado com especificações SDD+TDD completas, totalizando
24 CTs implementados e validados. A sinergia entre as especificações foi
calculada através de validação cruzada:
\begin{itemize}
\item SPEC-019 $\leftrightarrow$ SPEC-020: 0,85 (API Governance gerencia
      producers/consumers de streaming);
\item SPEC-019 $\leftrightarrow$ SPEC-021: 0,80 (Low-Code Platform expõe APIs
      governadas via Registry);
\item SPEC-020 $\leftrightarrow$ SPEC-021: 0,65 (Agentes low-code consomem
      streams tipados via Schema Registry).
\end{itemize}

\subsection{Artigo CORA-OpenCode}

O artigo \destaque{CORA-Eval: Um Framework de Benchmarking para Ciências
Exatas e da Natureza com Agentes Autônomos} \cite{laranjeira2026cora}
apresenta o CORA-Eval como framework de benchmark com:

\begin{itemize}
\item 150 tarefas em 10 dimensões $\times$ 4 níveis (Básico $\rightarrow$ Pesquisa);
\item Rastreador Python com persistência JSON;
\item Integração Cora V1-V7;
\item Q-Score UCB1 para seleção adaptativa;
\item Baseline CORA-Score 0,67;
\item CORA-V-Score ponderado por verificadores ativos.
\end{itemize}

O artigo foi submetido a periódico \qualisA{} na área de Ciência da Computação
e encontra-se em processo de revisão.

\subsection{Ensaio Qualis A1 (ensaio\_qualis\_a1.tex)}

O \destaque{Ensaio Qualis A1}, disponível no arquivo
\codigo{ensaio\_qualis\_a1.tex}, é um artigo de posicionamento que sintetiza
a trajetória de pesquisa do \opencodeeco{}. O ensaio aborda:

\begin{itemize}
\item A evolução de CLIs convencionais para ecossistemas cognitivos;
\item O papel da metacognição funcional em sistemas autônomos;
\item A integração de 600+ componentes em uma arquitetura coerente;
\item As lições aprendidas ao longo de 23 ciclos evolutivos (R1-R23);
\item As implicações para a engenharia de software do futuro.
\end{itemize}

\subsection{Artigo MIT/IA (artigo-mit-ia)}

O \destaque{Artigo MIT/IA} \cite{laranjeira2026opencode}, localizado no
diretório \codigo{artigo-mit-ia} (118 arquivos), foi submetido a periódico
internacional de alto impacto. O artigo aborda:

\begin{itemize}
\item Integração de motores de raciocínio formal (Z3, SymPy, miniKanren,
      Critical) em ecossistemas de agentes;
\item Validação experimental com 312 testes (312/312 PASS -- 100\%);
\item Arquitetura multiescalar com sincronização via Nexus;
\item Mecanismos de confiança e governança comportamental.
\end{itemize}

\subsection{Pipeline Completo: SEEKER $\rightarrow$ MASWOS $\rightarrow$ Correção $\rightarrow$ Validação $\rightarrow$ Publicação}

O pipeline de produção de artigos \qualisA{} segue o fluxo:

\begin{enumerate}
\item \destaque{SEEKER (Pesquisa)}: 10 agentes de pesquisa realizam busca
      sistemática em 10+ fontes acadêmicas, construindo uma base de
      conhecimento fundamentada \cite{opencode2026seeker};
\item \destaque{MASWOS (Escrita)}: 49 agentes especializados produzem o
      manuscrito conforme template \texttt{abntex2}{} e normas ABNT
      \cite{opencode2026maswos};
\item \destaque{Correção}: Iterative Correction Loop com 5 revisores,
      4 orientadores, 6 motores de correção;
\item \destaque{Validação}: PhD Auditor com Nash, Cohen, Bonferroni, Qualis,
      análise de sensibilidade;
\item \destaque{Publicação}: formatação final, geração de PDF, submissão ao
      periódico alvo.
\end{enumerate}

\begin{figure}[htb]
\centering
\begin{tikzpicture}[node distance=1.8cm, scale=0.5, every node/.style={scale=0.5},
    box/.style={draw, rounded corners, minimum width=2.5cm, minimum height=0.8cm, font=\small, align=center},
    arrow/.style={->, thick}]
\node[box, fill=blue!10] (seeker) {SEEKER\\10 agentes};
\node[box, fill=green!10, right=of seeker] (maswos) {MASWOS\\49 agentes};
\node[box, fill=orange!10, right=of maswos] (corr) {Correção\\15+ motores};
\node[box, fill=red!10, right=of corr] (valid) {Validação\\PhD Auditor};
\node[box, fill=purple!10, right=of valid] (pub) {Publicação\\Periódico A1};

\draw[arrow] (seeker) -- (maswos);
\draw[arrow] (maswos) -- (corr);
\draw[arrow] (corr) -- (valid);
\draw[arrow] (valid) -- (pub);
\end{tikzpicture}
\caption{Pipeline de produção de artigos \qualisA{}}
\label{fig:pipeline-artigos}
\end{figure}

\subsection{Estratégia de Submissão para Periódicos Qualis A1}

A estratégia de submissão segue as seguintes etapas:

\begin{enumerate}
\item \destaque{Seleção do periódico}: identificar periódicos \qualisA{} na
      área, analisar escopo editorial, fator de impacto, tempo de revisão;
\item \destaque{Template}: adequar o manuscrito ao template específico do
      periódico (cada periódico possui formatação própria);
\item \destaque{Carta de submissão}: redigir carta de submissão destacando a
      originalidade e relevância do trabalho;
\item \destaque{Peer review}: responder aos revisores de forma estruturada,
      utilizando as estratégias de defesa do Agent Forum;
\item \destaque{Revisão final}: incorporar as correções sugeridas pelos
      revisores e submeter a versão final;
\item \destaque{Acompanhamento}: monitorar o status da submissão e responder
      prontamente a solicitações adicionais.
\end{enumerate}

A Tabela~\ref{tab:estrategia-submissao} apresenta um cronograma típico de
submissão.

\begin{table}[htb]
\centering
    \small
\caption{Cronograma típico de submissão para periódico \qualisA{}}
\label{tab:estrategia-submissao}
    \resizebox{\textwidth}{!}{
\begin{tabular}{lll}
\toprule
\textbf{Etapa} & \textbf{Prazo} & \textbf{Ferramentas} \\
\midrule
Seleção do periódico & 2 semanas & SEEKER, Qualis A1 Auditor \\
Template & 1 semana & IMRAD Formatter \\
Carta de submissão & 2 dias & Agente 44 \\
Submissão inicial & 1 dia & Pipeline MASWOS \\
Resposta aos revisores & 4 semanas & Agent Forum \\
Revisão final & 2 semanas & Iterative Correction Loop \\
Publicação & Variável & Acompanhamento \\
\bottomrule
\end{tabular}
    }
\end{table}

\begin{exercicio}[Nivel PhD -- Pesquisa]
    Identifique 5 periódicos \qualisA{} na sua área de pesquisa. Para cada
    um, analise: escopo editorial, fator de impacto JCR, tempo médio de
    revisão e taxa de aceitação.
\end{exercicio}

\begin{exercicio}[Nivel PhD -- Escrita]
    Utilize o pipeline SEEKER $\rightarrow$ MASWOS para produzir um artigo
    curto (4-6 páginas) sobre um tema de sua escolha. Execute o Iterative
    Correction Loop até score $\geq$ 95 e submeta ao PhD Auditor.
\end{exercicio}

% =============================================================================
%  SEÇÃO 8.9: ROTEIRO DO NÍVEL ZERO AO PhD
%  Nível: Todos (~6 laudas)
% =============================================================================
\section{Roteiro do Nível Zero ao PhD}
\label{sec:roteiro-zero-phd}
\nivelzero

\paragraph{O mapa do tesouro do aprendiz a pesquisador} Esta seção apresenta um roteiro completo de aprendizado, organizado em marcos
de progressão que guiam o estudante do nível zero (nenhum conhecimento prévio)
ao nível PhD (capacidade de produzir pesquisa original e defendê-la perante
banca acadêmica).

\subsection{Roadmap Completo de Aprendizado}

A Figura~\ref{fig:roadmap} apresenta o roadmap visual de progressão.

\begin{figure}[htb]
\centering
\begin{tikzpicture}[
    node distance=1.2cm,
    milestone/.style={draw, rounded corners, minimum width=5cm, minimum height=0.8cm, font=\small\bfseries, align=center},
    level/.style={draw, rounded corners, minimum width=4.5cm, minimum height=0.6cm, font=\small, align=center},
    arrow/.style={->, thick}
]
\node[milestone, fill=red!10] (n0) {Nível 0 -- Fundamentos};
\node[level, below=0.3cm of n0] (n0c) {Lógica, Matemática básica, Raciocínio};

\node[milestone, fill=orange!10, below=0.8cm of n0c] (n1) {Nível Básico -- Programação};
\node[level, below=0.3cm of n1] (n1c) {Algoritmos, Python, Controle de versão};

\node[milestone, fill=yellow!10, below=0.8cm of n1c] (n2) {Nível Intermediário -- IA};
\node[level, below=0.3cm of n2] (n2c) {ML, Redes Neurais, Agentes simples};

\node[milestone, fill=green!10, below=0.8cm of n2c] (n3) {Nível Avançado -- Arquitetura};
\node[level, below=0.3cm of n3] (n3c) {Ecossistemas, MCPs, Skills, Agentes};

\node[milestone, fill=blue!10, below=0.8cm of n3c] (n4) {Nível PhD -- Pesquisa};
\node[level, below=0.3cm of n4] (n4c) {Metacognição, Trust, Validação científica};

\draw[arrow] (n0c.south) -- (n1.north);
\draw[arrow] (n1c.south) -- (n2.north);
\draw[arrow] (n2c.south) -- (n3.north);
\draw[arrow] (n3c.south) -- (n4.north);
\end{tikzpicture}
\caption{Roadmap de aprendizado do nível zero ao PhD}
\label{fig:roadmap}
\end{figure}

A Tabela~\ref{tab:marcos-progressao} detalha os marcos de progressão,
incluindo os pré-requisitos, competências adquiridas e certificações
associadas.

\begin{table}[htb]
\centering
    \small
\caption{Marcos de progressão do nível zero ao PhD}
\label{tab:marcos-progressao}
    \resizebox{\textwidth}{!}{
\begin{tabular}{p{2cm}p{3cm}p{4cm}p{3cm}}
\toprule
\textbf{Nível} & \textbf{Pré-requisitos} & \textbf{Competências adquiridas} & \textbf{Capítulos} \\
\midrule
0 -- Zero & Nenhum & Lógica, conjuntos, funções, raciocínio formal & 1 \\
Básico & Nível 0 & Python, algoritmos, estruturas de dados, Git & 1 \\
Intermediário & Básico & ML, redes neurais, LLMs, prompts, agentes & 2 \\
Avançado & Intermediário & Arquitetura 3 camadas, MCPs, skills, plugins & 3, 4, 5 \\
PhD & Avançado & Metacognição, trust, economia, validação, defesa & 6, 7, 8 \\
\bottomrule
\end{tabular}
    }
\end{table}

\subsection{[Detalhamento] Marcos de Progressão}

\subsubsection{Nível 0: Fundamentos Matemáticos e Lógica}

O estudante inicia sua jornada dominando os fundamentos matemáticos que
sustentam toda a ciência da computação \cite{russell2020inteligencia}:

\begin{itemize}
\item Lógica proposicional e de primeira ordem;
\item Teoria dos conjuntos e funções;
\item Álgebra linear (vetores, matrizes, transformações);
\item Cálculo diferencial e integral;
\item Probabilidade e estatística descritiva.
\end{itemize}

\textbf{Carga horária estimada}: 80 horas.
\textbf{Projetos práticos}: implementar um provador de teoremas simples em
Python; calcular métricas estatísticas de um dataset real.

\subsubsection{Nível Básico: Programação e Algoritmos}

Com os fundamentos matemáticos estabelecidos, o estudante avança para a
programação \cite{cormen2009algorithms}:

\begin{itemize}
\item Python: sintaxe, estruturas de dados, orientação a objetos;
\item Algoritmos: busca, ordenação, grafos, programação dinâmica;
\item Controle de versão com Git/GitHub;
\item Linha de comando e scripting.
\end{itemize}

\textbf{Carga horária estimada}: 120 horas.
\textbf{Projetos práticos}: implementar um interpretador de expressões
lógicas; criar um sistema de busca textual simples.

\subsubsection{Nível Intermediário: IA e Agentes}

O estudante mergulha na inteligência artificial e nos sistemas multiagentes
\cite{russell2020inteligencia, wooldridge2009multiagent}:

\begin{itemize}
\item Aprendizado de máquina supervisionado e não supervisionado;
\item Redes neurais e deep learning;
\item Modelos de linguagem (LLMs) e engenharia de prompts;
\item Fundamentos de agentes inteligentes.
\end{itemize}

\textbf{Carga horária estimada}: 200 horas.
\textbf{Projetos práticos}: treinar um classificador de texto; implementar
um agente reativo simples.

\subsubsection{Nível Avançado: Arquitetura de Ecossistemas}

O estudante aprende a arquitetura e engenharia de ecossistemas cognitivos
\cite{opencode2026ecosystem}:

\begin{itemize}
\item Arquitetura três camadas (MCP $\rightarrow$ Skill $\rightarrow$ Agent);
\item Metodologia SDD+TDD para especificação e teste;
\item Scanner Pipeline (Noológico, Teleológico, Evolutivo);
\item Trust Engine e Behavioral Gate;
\item Token Economy e Economia de Agentes.
\end{itemize}

\textbf{Carga horária estimada}: 300 horas.
\textbf{Projetos práticos}: implementar uma skill personalizada; configurar
um novo MCP; criar um agente especializado.

\subsubsection{Nível PhD: Metacognição, Trust, Validação Científica}

O estudante atinge o nível máximo de proficiência, capaz de realizar pesquisa
original e defendê-la \cite{popper2005logica, opencode2026spec036}:

\begin{itemize}
\item Metacognição funcional (SPEC-036);
\item Structural Noise Scanner (SPEC-037);
\item Trust Engine (SPEC-038);
\item PhD Auditor: Nash, Cohen, Bonferroni, Qualis;
\item Produção de artigos \qualisA{};
\item Defesa perante banca acadêmica.
\end{itemize}

\textbf{Carga horária estimada}: 400 horas.
\textbf{Projetos práticos}: produzir um artigo \qualisA{} completo; simular
uma banca de defesa; validar o ecossistema com 312 testes.

\subsection{Como Usar Este Livro como Guia Autodidata}

Este livro foi estruturado para permitir o estudo autodidata progressivo.
Recomenda-se:

\begin{enumerate}
\item \destaque{Leitura sequencial}: seguir os capítulos na ordem apresentada,
      do Capítulo 1 ao Capítulo 8;
\item \destaque{Prática constante}: executar todos os exemplos de código e
      resolver todos os exercícios;
\item \destaque{Projetos acumulativos}: aplicar os conceitos de cada capítulo
      em um projeto pessoal que cresce ao longo do livro;
\item \destaque{Comunidade}: participar da comunidade do \opencodeeco{} no
      GitHub, contribuindo com issues, pull requests e discussões.
\end{enumerate}

\subsection{Próximos Passos: Pós-Doutorado e Pesquisa Avançada}

Para o pesquisador que completa o nível PhD, as próximas fronteiras incluem:

\begin{itemize}
\item \destaque{Pesquisa em consciência artificial}: explorar o Self-Model
      N0-N3 e suas implicações filosóficas;
\item \destaque{Economia de agentes em larga escala}: expandir a Token Economy
      para milhares de agentes;
\item \destaque{Validação clínica}: aplicar o pipeline de validação a domínios
      como medicina e biologia computacional;
\item \destaque{Governança descentralizada}: implementar DAOs (Organizações
      Autônomas Descentralizadas) para governança do ecossistema;
\item \destaque{Publicação internacional}: submeter os resultados a periódicos
      internacionais de alto impacto.
\end{itemize}

\begin{exercicio}[Nivel 0 -- Autoavaliação]
    Autoavalie seu nível atual em cada um dos 5 marcos de progressão
    (0 a PhD). Para cada nível não atingido, liste as competências que
    você precisa desenvolver e estime a carga horária necessária.
\end{exercicio}

\begin{exercicio}[Nivel Básico -- Planejamento]
    Crie um plano de estudos personalizado de 6 meses para avançar do seu
    nível atual ao próximo nível. Inclua metas semanais, recursos de estudo
    e projetos práticos.
\end{exercicio}

\begin{exercicio}[Nivel Intermediário -- Projeto]
    Selecione um dos projetos práticos sugeridos para seu nível atual e
    implemente-o utilizando o \opencodeeco{}. Documente o processo e os
    resultados em formato de artigo curto.
\end{exercicio}

\begin{exercicio}[Nivel Avançado -- Extensão]
    Identifique um gap no \opencodeeco{} que você possa endereçar com uma
    nova skill, MCP ou plugin. Implemente a extensão, documente-a em formato
    SPEC e submeta como pull request no repositório do ecossistema.
\end{exercicio}

\begin{exercicio}[Nivel PhD -- Publicação]
    Produza um artigo completo (8-12 páginas) sobre sua pesquisa utilizando
    o pipeline SEEKER-MASWOS-Correção-Validação do \opencodeeco{}. Submeta
    o artigo a um periódico \qualisA{} da sua área.
\end{exercicio}

% =============================================================================
%  SEÇÃO 8.10: CONCLUSÃO E PERSPECTIVAS FUTURAS
%  Nível: PhD (~6 laudas)
% =============================================================================
\section{Conclusão e Perspectivas Futuras}
\label{sec:conclusao-perspectivas}
\nivelphd

\paragraph{Fechando ciclos e abrindo novos horizontes} Este capítulo final do livro sintetiza a jornada completa de pesquisa e
desenvolvimento do \opencodeeco{}, desde sua concepção como uma CLI
convencional até sua maturidade como um ecossistema cognitivo com 600+
componentes integrados.

\subsection{Síntese da Jornada: do R1 ao R23}

A Tabela~\ref{tab:ciclos-resumo} apresenta um resumo dos 23 ciclos evolutivos,
destacando as principais contribuições de cada ciclo.

\begin{table}[htb]
\centering
    \small
\caption{Resumo dos ciclos evolutivos R1-R23}
\label{tab:ciclos-resumo}
    \resizebox{\textwidth}{!}{
\begin{tabular}{clc}
\toprule
\textbf{Ciclo} & \textbf{Contribuição principal} & \textbf{Score} \\
\midrule
R1 & Cross-validation quantitativa, dados World Bank & 85 \\
R2 & Pipeline de artigos acadêmicos & 90 \\
R3 & TSAC, Sci-Hub, cross-validation refinada & 92 \\
R4-R5 & Ciclo de correção iterativa, CJK detection & 95-98 \\
R6-R7 & Editais-br, cache versionado & 92-94 \\
R8-R9 & SDD+TDD acadêmico, AutoEvolve LaTeX & 94-96 \\
R10 & Menu adaptativo, plugin system & 96 \\
R11 & CORA-Eval Benchmark & 97 \\
R12 & Science Skills + MCP Expansion & 98 \\
R13 & Reasoning Engines (Z3, SymPy, Kanren, Critical) & 96 \\
R14-R16 & Expansão (227 skills, 128 agentes, 46 MCPs) & 97-98 \\
R17 & Gartner Hype Cycle 2026, 3 gaps estratégicos & 99 \\
R18-R18b & Token Economy, Agent Economics, Audit & 99 \\
R19 & MCSP + Scanner Ecosystem (76 CTs) & 99 \\
R20 & Composição Unitária do Conhecimento & 100 \\
R21 & Metacognição + Self-Evolution (SPEC-036) & 100 \\
R22 & Structural Noise Scanner + N3 Completo (SPEC-037) & 100 \\
R23 & Trust Engine + Behavioral Autonomy (SPEC-038) & 100 \\
\bottomrule
\end{tabular}
    }
\end{table}

A evolução de R1 (score 85) a R23 (score 100) demonstra o amadurecimento
contínuo do ecossistema, impulsionado pelos mecanismos de autoavaliação e
evolução autônoma \cite{opencode2026ecosystem}.

\subsection{Contribuições Originais da Pesquisa}

As contribuições originais desta pesquisa, sistematicamente validadas ao
longo dos 23 ciclos evolutivos, incluem:

\begin{enumerate}
\item \destaque{Arquitetura três camadas (MCP-Skill-Agent)}: um padrão
      arquitetural que separa claramente as responsabilidades de
      infraestrutura, habilidade e agência, permitindo composição flexível
      e reuso de componentes;
\item \destaque{Scanner Pipeline (Noológico-Teleológico-Evolutivo-
      Refinamento-MCSP)}: um pipeline de 5 scanners que analisam o código
      em múltiplas dimensões --- estado atual, estado futuro, trajetórias
      evolutivas, refinamento e solução de capacidade mínima;
\item \destaque{Metacognição Funcional (SPEC-036)}: implementação de
      auto-monitoramento, forecasting, introspecção de fonte, boundary
      self/other e análise causal (Granger + Bayes) em sistemas de
      engenharia de software;
\item \destaque{Trust Engine (SPEC-038)}: sistema de confiança com
      Behavioral Gate preventivo, Natural Forgetting (Atkinson-Shiffrin)
      e TrustScorer blend 70/30, garantindo segurança em operação autônoma;
\item \destaque{Token Economy (SPEC-022/023/024)}: sistema de incentivos
      econômicos para agentes, com ledger congelado, fee market dinâmico,
      staking com lock de 7 dias e slashing;
\item \destaque{PhD Auditor e AUTO\_SCORE\_QUALIS}: sistema de validação
      científica que aplica rigor estatístico (Nash, Cohen, Bonferroni)
      à avaliação acadêmica;
\item \destaque{Protocolo de Anonimato}: sistema de detecção e remoção de
      identificadores diretos e indiretos para avaliação cega.
\end{enumerate}

\subsection{Limitações e Trabalhos Futuros}

Como toda pesquisa científica, este trabalho possui limitações que apontam
direções para investigações futuras:

\begin{enumerate}
\item \destaque{Dependência de LLMs externos}: o ecossistema depende de APIs
      de modelos de linguagem de grande escala (DeepSeek, GPT-4), o que
      introduz dependência de terceiros e custos operacionais. \textbf{Trabalho
      futuro}: implementar um modelo local fine-tunado específico para
      engenharia de software;
\item \destaque{Escalabilidade horizontal}: a arquitetura atual foi testada
      com dezenas de agentes simultâneos. \textbf{Trabalho futuro}: validar
      a arquitetura com centenas ou milhares de agentes em paralelo;
\item \destaque{Generalização para outros domínios}: o ecossistema foi
      desenvolvido e validado no contexto de engenharia de software.
      \textbf{Trabalho futuro}: adaptar e validar a arquitetura para
      domínios como medicina, direito e educação;
\item \destaque{Segurança formal}: o Behavioral Gate é preventivo, mas não
      oferece garantias formais de segurança. \textbf{Trabalho futuro}:
      integrar verificação formal (Z3, Lean) ao pipeline de autorização;
\item \destaque{Consciência artificial}: o Self-Model N0-N3 é uma
      aproximação funcional, não uma implementação de consciência.
      \textbf{Trabalho futuro}: explorar modelos de consciência baseados
      em teoria da informação integrada (IIT) \cite{tononi2008consciousness}.
\end{enumerate}

\subsection{O Futuro dos Ecossistemas Cognitivos}

O campo de ecossistemas cognitivos artificiais está em rápida evolução.
As tendências que moldarão seu futuro incluem:

\begin{itemize}
\item \destaque{Agentes multimodalidade}: integração de visão computacional,
      processamento de áudio e interação física em agentes cognitivos;
\item \destaque{Economia descentralizada de agentes}: mercados autônomos onde
      agentes negociam serviços, dados e recursos computacionais;
\item \destaque{Governança algorítmica}: sistemas de governança baseados em
      contratos inteligentes e DAOs para coordenação de agentes;
\item \destaque{Aprendizado contínuo}: agentes que aprendem continuamente com
      a experiência, sem esquecimento catastrófico;
\item \destaque{Alinhamento verificável}: garantias formais de que o
      comportamento dos agentes está alinhado com valores e objetivos humanos;
\item \destaque{Ecossistemas auto-sustentáveis}: sistemas capazes de manter
      sua própria operação, evolução e governança sem intervenção humana.
\end{itemize}

O Gartner Hype Cycle 2026 \cite{gartner2026hype} posiciona tecnologias como
\textit{Agent AI}, \textit{Autonomous Systems} e \textit{AI Governance} no
pico de expectativas infladas, sugerindo que 2-5 anos serão necessários para
que atinjam maturidade produtiva. O \opencodeeco{}, com sua arquitetura
madura e validação científica robusta, está posicionado na vanguarda desta
onda tecnológica.

\subsection{Chamado à Ação: Contribua para o OpenCode Ecosystem}

O \opencodeeco{} é um projeto de código aberto que vive das contribuições da
comunidade. Convidamos o leitor a:

\begin{itemize}
\item \destaque{Experimentar}: instalar o ecossistema, explorar seus comandos
      e agentes, e utilizá-lo em projetos reais;
\item \destaque{Contribuir}: relatar bugs, sugerir funcionalidades, submeter
      pull requests e escrever documentação;
\item \destaque{Pesquisar}: utilizar o ecossistema como plataforma de pesquisa
      para experimentos em engenharia de software, IA e sistemas multiagentes;
\item \destaque{Ensinar}: utilizar este livro como material didático em cursos
      de graduação e pós-graduação;
\item \destaque{Publicar}: submeter artigos sobre suas experiências e
      extensões do ecossistema a periódicos e conferências.
\end{itemize}

O repositório do ecossistema está disponível em:
\url{https://github.com/marceloclaro/opencode-ecosystem}

A documentação completa e os guias de contribuição estão em:
\url{https://github.com/MarceloClaro/OpenCode_Ecosystem}

\begin{exercicio}[Nivel PhD -- Meta-reflexão]
    Reflita sobre sua jornada de aprendizado ao longo deste livro. Escreva um
    texto de 2-3 páginas respondendo: (a) Qual foi o conceito mais desafiador?
    (b) Qual foi a descoberta mais surpreendente? (c) Como você pretende
    aplicar este conhecimento na prática?
\end{exercicio}

\begin{exercicio}[Nivel PhD -- Contribuição]
    Identifique uma melhoria ou extensão para o \opencodeeco{} que você
    gostaria de implementar. Documente-a no formato SPEC (Spec-Driven
    Development) e submeta como issue no repositório do ecossistema.
\end{exercicio}

\begin{exercicio}[Nivel PhD -- Projeto Final]
    Este é o exercício culminante do livro: produza um artigo científico
    completo (8-12 páginas) sobre um tema de sua escolha, utilizando todo
    o pipeline do \opencodeeco{}:
    \begin{enumerate}
    \item SEEKER para pesquisa bibliográfica (mínimo 20 referências);
    \item MASWOS para escrita do manuscrito;
    \item Iterative Correction Loop até score $\geq$ 95;
    \item PhD Auditor para validação científica;
    \item Submissão a periódico ou conferência da área.
    \end{enumerate}
\end{exercicio}

\section*{Referências do Capítulo}

\begin{itemize}
\item Para metodologia de pesquisa em engenharia de software:
      \cite{sommerville2011engenharia} e \cite{sommerville2011engenharia};
\item Para normas ABNT de trabalhos acadêmicos:
      \cite{abnt2023nbr14724, abnt2023nbr6023, abnt2023nbr10520};
\item Para sistemas multiagentes e teoria dos jogos:
      \cite{wooldridge2009multiagent, weiss2013multiagent, nash1950equilibrium};
\item Para estatística e rigor científico:
      \cite{wasserman2004all, cohen2018estatistica, bonferroni1936teoria};
\item Para filosofia da ciência:
      \cite{popper2005logica, kuhn1998estrutura};
\item Para o ecossistema OpenCode:
      \cite{opencode2026ecosystem, opencode2026spec036, opencode2026spec038,
      opencode2026maswos, opencode2026seeker, opencode2026cora};
\item Para publicações e artigos:
      \cite{laranjeira2026opencode, laranjeira2026gartner,
      laranjeira2026cora, laranjeira2026trust};
\item Para o Gartner Hype Cycle:
      \cite{gartner2026hype};
\item Para ética em IA:
      \cite{bender2021dangers, gebru2021datasheets}.
\end{itemize}

\begin{observacao}
Todos os exemplos de código, scripts e pipelines descritos neste capítulo
estão disponíveis no repositório do \opencodeeco{} sob o diretório
\codigo{examples/capitulo8/}. O leitor é incentivado a executá-los,
modificá-los e utilizá-los como base para seus próprios projetos acadêmicos.
Para instalar o ecossistema, consulte a documentação oficial em
\url{https://github.com/marceloclaro/opencode-ecosystem}.
\end{observacao}

% =============================================================================
%  FIM DO CAPÍTULO 8
